% Einheit 3 von 5 – Winkelsummen, Dreiecke und Vierecke (bank/winkel-dreiecke/e3.jsonl, zusammenbau v0.1)
%% TODO Zweigzeile Teil 1: was hier gelernt wird (Satz in Schülersprache aus dem Einheitstitel) – Einheit 3
\einheitenkopf[e3]{Einheit 3 von 5 · Winkelsummen, Dreiecke und Vierecke}
\zweigzeile{ab Kl. 5 · P10 oft}
\setcounter{aufgabe}{15}

%% TODO Titel: Ich-kann-Satz fehlt in der Bank (Platzhalter: Kettenname) – Nr. 16
%% TODO Anweisung: ein Satz über der ersten Teilaufgabe fehlt in der Bank – Nr. 16
\begin{aufgabe}{Welches Teildreieck?}
\begin{teile}
\teil Im Dreieck ABC ist CD die Höhe auf AB. Gesucht ist der Winkel bei C zwischen AC und CD. Fahre in einer Skizze das Dreieck nach, in dem er liegt, und kreuze an: \\ \kreuz{Dreieck ADC} \\ \kreuz{Dreieck DBC} \\ \kreuz{Dreieck ABC}
\end{teile}
\end{aufgabe}

%% TODO Titel: Ich-kann-Satz fehlt in der Bank (Platzhalter: Kettenname) – Nr. 17
%% TODO Anweisung: ein Satz über der ersten Teilaufgabe fehlt in der Bank – Nr. 17
\begin{aufgabe}{Winkelsumme}
\begin{teile}
\teil Markiere gleich lange Seiten oder den rechten Winkel und kreuze die Dreiecksart an. \\ \kreuz{spitzwinklig} \\ \kreuz{rechtwinklig} \\ \kreuz{stumpfwinklig}
\end{teile}
%% TODO Darstellung für schwach fehlt in der Bank (winkel-dreiecke-e3-k2-s1-v1; Abschnitt „Für schwache Schüler“)
\swz{Im Dreieck ABC ist $\alpha = 52^\circ$ und $\beta = 71^\circ$. Wie groß ist γ?}{}
%% TODO Darstellung für schwach fehlt in der Bank (winkel-dreiecke-e3-k2-s1-v2; Abschnitt „Für schwache Schüler“)
\swz{Im Dreieck ABC ist $\alpha = 36^\circ$ und $\beta = 88^\circ$. Wie groß ist γ?}{}
%% TODO Darstellung für schwach fehlt in der Bank (winkel-dreiecke-e3-k2-s1-v3; Abschnitt „Für schwache Schüler“)
\swz{Im Dreieck ABC ist $\alpha = 29^\circ$ und $\beta = 103^\circ$. Wie groß ist γ?}{}
%% TODO Darstellung für schwach fehlt in der Bank (winkel-dreiecke-e3-k2-s1-v4; Abschnitt „Für schwache Schüler“)
\swz{Im Dreieck ABC ist $\alpha = 64^\circ$ und $\beta = 59^\circ$. Wie groß ist γ?}{}
%% TODO Darstellung für schwach fehlt in der Bank (winkel-dreiecke-e3-k2-s1-v5; Abschnitt „Für schwache Schüler“)
\swz{Im Dreieck ABC ist $\alpha = 43^\circ$ und $\beta = 95^\circ$. Wie groß ist γ?}{}
\swfrage{Was bleibt gleich, was ändert sich?}
%% TODO Darstellung für schwach fehlt in der Bank (winkel-dreiecke-e3-k2-s2-v1; Abschnitt „Für schwache Schüler“)
\swz{Ein dreieckiges Grundstück hat an zwei Ecken die Winkel $51{,}4^\circ$ und $76{,}9^\circ$. Wie groß ist der Winkel an der dritten Ecke? \hfill (P10 2017 OS)}{}
%% TODO Darstellung für schwach fehlt in der Bank (winkel-dreiecke-e3-k2-s3-v1; Abschnitt „Für schwache Schüler“)
\swz{Ein gleichschenkliges Dreieck hat zwei Basiswinkel von je $67^\circ$. Wie groß ist der Winkel an der Spitze?}{}
%% TODO Darstellung für schwach fehlt in der Bank (winkel-dreiecke-e3-k2-s4-v1; Abschnitt „Für schwache Schüler“)
\swz{Ein Dreieck hat drei gleich lange Seiten. Wie groß ist jeder seiner Winkel?}{}
%% TODO Darstellung für schwach fehlt in der Bank (winkel-dreiecke-e3-k2-s5-v1; Abschnitt „Für schwache Schüler“)
\swz{Ein Dreieck hat einen rechten Winkel und einen Winkel von $38^\circ$. Wie groß ist der dritte Winkel?}{}
%% TODO Darstellung für schwach fehlt in der Bank (winkel-dreiecke-e3-k2-s6-v1; Abschnitt „Für schwache Schüler“)
\swz{Ein Viereck hat die Winkel $84^\circ$, $97^\circ$ und $103^\circ$. Wie groß ist der vierte Winkel?}{}
%% TODO Darstellung für schwach fehlt in der Bank (winkel-dreiecke-e3-k2-s7-v1; Abschnitt „Für schwache Schüler“)
\swz{Im Giebeldreieck ABC eines Hauses steht die Stütze CD senkrecht auf dem Balken AB. An der Ecke A ist $\alpha = 27{,}6^\circ$. Wie groß ist der Winkel $\gamma_1$ zwischen AC und der Stütze CD? \hfill (P10 2022 OS)}{}
%% TODO Darstellung für schwach fehlt in der Bank (winkel-dreiecke-e3-k2-s8-v1; Abschnitt „Für schwache Schüler“)
\swz{Im Dreieck ABC ist $\alpha = \beta = 64^\circ$ und AC = 5 cm. Welche Seite ist genauso lang wie AC, und wie lang ist sie?}{}
%% TODO Darstellung für schwach fehlt in der Bank (winkel-dreiecke-e3-k2-s9-v1; Abschnitt „Für schwache Schüler“)
\swz[3]{Im Dreieck ABC sind die Winkel bei A und bei B je halb so groß wie ein rechter Winkel. Begründe, dass der Winkel bei C ein rechter ist.}{}
\begin{teile}
\teil Ergänze richtig: „In jedem Drachenviereck …“ \hfill (P10 2026 FOR) \\ \kreuz{sind alle Winkel gleich groß} \\ \kreuz{gibt es zwei Paare gleich langer Nachbarseiten} \\ \kreuz{sind alle Seiten gleich lang}
\end{teile}
\swz[3]{Im Dreieck ABC ist $\alpha = \beta = 66^\circ$ und AC = 7 cm. Gib die Länge von BC und die Größe von γ an. \hfill (P10 2023 OS)}{}
\end{aufgabe}

%% TODO Titel: Ich-kann-Satz fehlt in der Bank (Platzhalter: Kettenname) – Nr. 18
%% TODO Anweisung: ein Satz über der ersten Teilaufgabe fehlt in der Bank – Nr. 18
\begin{aufgabe}{Innenwinkelsumme von Fünf- und Sechseck über Dreiecke}
%% TODO Darstellung für schwach fehlt in der Bank (winkel-dreiecke-e3-k3-s1-v1; Abschnitt „Für schwache Schüler“)
\swz{Zerlege ein Fünfeck durch Diagonalen von einer Ecke aus in Dreiecke. Wie groß ist die Summe seiner Innenwinkel?}{}
\end{aufgabe}

%% TODO Titel: Ich-kann-Satz fehlt in der Bank (Platzhalter: Kettenname) – Nr. 19
%% TODO Anweisung: ein Satz über der ersten Teilaufgabe fehlt in der Bank – Nr. 19
\begin{aufgabe}{Winkelsumme – Fehler finden, Begründen, Anwendung}
\swz[3]{Ein Viereck hat die Winkel $52^\circ$, $71^\circ$ und $93^\circ$. Ben rechnet den vierten Winkel: \rechnung{52^\circ + 71^\circ + 93^\circ &= 216^\circ \\ 216^\circ - 180^\circ &= 36^\circ} Finde den Fehler und rechne richtig.}{}
\swfrage{Zeichne durch C die Parallele zu AB. Begründe mit Wechselwinkeln, dass α, β und γ zusammen einen gestreckten Winkel ergeben.}
\swz[3]{Ein Satteldach ist im Querschnitt ein gleichschenkliges Dreieck; beide Dachflächen sind um $38^\circ$ gegen die Waagerechte geneigt. Wie groß ist der Winkel am First?}{}
\end{aufgabe}

\uebersichtskasten{Dreieck: Die drei Innenwinkel ergeben zusammen 180°. Viereck: 360°. \\ α = 61°, β = 77°: γ = 180° − 61° − 77° = 42° \quad Viereck mit 90°, 90°, 112°: vierter Winkel 68° \\ Gleichschenklig: zwei gleich lange Schenkel, die beiden Basiswinkel sind gleich groß – und umgekehrt: gleiche Basiswinkel bedeuten gleich lange Schenkel. \quad Basiswinkel 74°: Spitze 180° − 2 · 74° = 32° \\ Gleichseitig: alle Winkel 60°. Rechtwinklig: die zwei spitzen Winkel ergeben zusammen 90°. \quad rechtwinklig mit 33°: dritter Winkel 57° \\ Vierecke: Trapez – ein Paar paralleler Seiten. Parallelogramm – zwei Paare, Gegenwinkel gleich. Raute – Parallelogramm mit vier gleich langen Seiten. Rechteck – vier rechte Winkel. Quadrat – beides. Drachen – zwei Paare gleich langer Nachbarseiten, Diagonalen senkrecht.}
